\section{Concrete Lean, Forge, Leant, and Djex integration}
\label{sec:integration}

\subsection{Do not invent a new outer controller}
The reviewed Forge runtime sketch already has obligations, scope keys,
conditional candidates, and replies named \code{candidate}, \code{conditional},
\code{progress}, and \code{noResult}. Those remain the right entry points for
positive proof construction. They are sketches rather than an implemented
\code{forge} runtime \cite{forge-runtime}.

The new positive worker should return a \code{ConditionalCandidate} whose
premises are the local polynomial identities, initial equalities, and a
source-to-machine semantic bridge. Its constructor is the appropriate
induction theorem. After those premises have Lean proofs, the ordinary Forge
acceptance path may install the result. A Python Boolean is never submitted as
that proof.

The Kripke lane should return a separate \emph{diagnostic artifact}. The
current reply sketch does not distinguish object-logic nonderivability from a
source-goal refutation. Add a nominally separate artifact carrying its exact
object formula, atom-origin map, and a proof about a selected derivation type.
Do not encode it as a candidate for $\neg G$.

\subsection{The smallest useful source-to-machine reifier}
The first implementation should accept elaborated, monomorphic folds with
fixed products of rational registers and explicit list element types. It can
be built through six checkable transformations:
\begin{enumerate}
\item Locate an equality between supported folds or their fixed polynomial
      output projections. Obtain the actual elaborated step functions and
      initial states, not a reparsed display string.
\item Normalize tuple projections and let-bindings using proved equations;
      distinguish old-state register expressions from newly computed values.
\item Reify literals, addition, multiplication, subtraction over $\Q$, and
      natural powers into a typed polynomial AST. Refuse all other arithmetic
      operations in the first version.
\item Reify complete finite constructor matches as control edges. Every
      constructor branch must have a corresponding transition; do not
      interpret a partial pattern match as a total transition system.
\item Build the product machine with one shared payload tuple per paired step.
      Prove the equality between each original step and evaluation of its
      reified polynomial update.
\item Instantiate the fold/run correspondence and translate output equality
      into the demanded polynomial difference. Keep this bridge among the
      candidate's explicit proof obligations.
\end{enumerate}

These operations can initially be implemented in ordinary \code{MetaM} using
Lean expression construction and standard elaboration. No undocumented
incremental \code{grind} hook is required for this worker. The arithmetic
identities can be discharged in separate local goals, and the completed
invariant can then be supplied to \code{grind} as an ordinary lemma.

This makes the contribution orthogonal to local saturation. The documented
\code{grind} engine remains useful for closing local consequences; it is not
being misdescribed as lacking congruence closure, arithmetic, or E-matching
\cite{grind}. The added capability is choosing and proving a finite invariant
that expresses an unbounded execution property.

\subsection{The initial Lean proof path: source reconstruction}
For each accepted certificate, emit the following obligations over the original
register and input binders:
\[
 P_{q_0}(x_0)=0,\qquad G_q=C_qP_q,\qquad
 P_{q'}\circ F_e=M_eP_q.
\]
Use \code{norm_num} for closed rational initial calculations and
\code{ring} or \code{ring_nf} for the universal polynomial identities.
Then use an ordinary induction theorem over the machine's reachability or
input-list execution.

The package's \code{emit_lean.py} performs the local source-emission part.
It emits 222 named identity obligations from the 37 positive certificates into
\code{lean/GeneratedIdentities.lean}. The file includes \code{\#print axioms}
commands for later auditing. It does not emit or prove a bridge from arbitrary
original Lean source, and no compiler was available here to validate those
identities. The generated count is a count of \emph{local obligations}, not
222 proved user theorems.

\code{lean/Reachability.lean} supplies complete, but uncompiled, proof text for
three generic induction bridges: reachable-state invariant preservation,
fold invariant preservation, and relational fold preservation. A hand-written
delayed-accumulator reconstruction appears in \code{lean/WorkedFold.lean}.
These distinguish a concrete proof path from an abstract promise to ``replay in
Lean.'' Their uncompiled status remains a material limitation.

\subsection{A reflected checker is the next optimization, not a prerequisite}
For larger instances, prove one theorem of the schematic form
\begin{lstlisting}
checkClosure problem certificate = true ->
  forall state, Reachable problem state -> GoalsHold problem state
\end{lstlisting}
inside Lean. The data types need a semantic interpretation into the same
rational polynomial ring used by the source bridge. Soundness factors into
canonical decoding, correctness of sparse operations, the evaluated identity
lemma, and the reachability induction theorem.

Mathlib's multivariate polynomial evaluation infrastructure supplies relevant
semantics and homomorphism lemmas \cite{mvpoly}. Its representation need not
be the fastest executable representation. A compact sparse AST can carry its
own denotation into \code{MvPolynomial}, with a proved correctness theorem for
normalization. The external worker can retain any faster internal data structure
because Lean checks only the emitted certificate.

Clearing denominators row by row can reduce the rational checker to integer
polynomial identities. Record each positive denominator and multiply the
whole corresponding equation, rather than round coefficients. Prove the
nonzero scalar cancellation once. This is a deterministic representation
optimization; it does not alter the certificate's mathematical claim.

\subsection{Universal behavioral assertions in Leant}
At the reviewed Leant snapshot, the named \code{:synth f : TYPE where PROP}
route checks the exact proposed candidate and attempts positive and negative
proofs with \code{decide} and bounded simplification. A quantified assertion can
remain inconclusive when those methods do not close it
\cite{leant-behavior}.

Extend that assertion checker with a new dispatch rule, not a new inhabitant
search engine. When the asserted property is a supported universal fold
equality, feed the exact candidate/reference pair to the product-machine
reifier and observable-closure worker. A fully reconstructed Lean proof makes
the behavioral assertion pass. A fully reconstructed concrete source execution
can falsify it. A reifier refusal, search cutoff, or unproved semantic bridge
leaves the assertion inconclusive.

The candidate's original exact type and source identity remain unchanged.
A successful machine certificate for one candidate cannot be used for another
candidate with the same displayed type. This reuses the existing source-owned
candidate discipline rather than claiming it as an invention.

\subsection{Source-level quantifiers and arithmetic domains}
A closed rational initial vector is the prototype's direct input. Universally
quantified rational initial parameters can use the special initialization-edge
construction in Section~\ref{sec:model}. More general parameter telescopes
require a source bridge that proves this encoding preserves the quantifier
order and variable dependencies.

Positive rational polynomial identities also evaluate in a commutative
$\Q$-algebra. This gives a route to real-valued versions of the same equalities,
but the initial implementation must not silently cast an arbitrary source
semantics into that setting. A rational trace only becomes a source-level
counterexample after its values are well typed and its operations are proved
to match the source operations. In particular, negative rational states do not
automatically provide counterexamples to natural-number programs.

\subsection{A precise Kripke-to-Lean target}
Choose an IPC syntax and derivation system already formalized in Lean. Define a
Boolean evaluator for finite models and prove that it agrees with the library's
semantic forcing relation. The desired output theorem is schematically
\begin{lstlisting}
checkCountermodel Gamma goal certificate = true ->
  Not (Nonempty (IPC.Derivation Gamma goal))
\end{lstlisting}
The inner \code{IPC.Derivation} is an object-language derivation type; the
statement is not \code{Not (denote goal)}.

An object sequent extracted from Djex's formula structure needs a faithful map
of constructor-tagged disjunctions and conjunction arities, and an exact atom
identity table. A successful object-model check does not prove the separate
translation-completeness theorem for opaque quantified source types or
class dictionaries. Until that theorem is available, report only the named
object-logic result. Classical fallback may still prove the user's source
formula.
